% Generated draft: requires mathematical and source review.
\documentclass[11pt]{article}
\usepackage[margin=25mm]{geometry}
\usepackage{fontspec}
\setmainfont{Latin Modern Roman}
\usepackage{amsmath,amssymb,mathtools,unicode-math}
\setmathfont{Latin Modern Math}
\usepackage{enumitem,needspace,xurl,hyperref}
\setlength{\parindent}{0pt}
\setlength{\parskip}{5pt}
\setlength{\emergencystretch}{3em}
\newcommand{\N}{\mathbb{N}}
\newcommand{\Z}{\mathbb{Z}}
\newcommand{\Q}{\mathbb{Q}}
\newcommand{\R}{\mathbb{R}}
\newcommand{\C}{\mathbb{C}}
\newcommand{\F}{\mathbb{F}}
\newcommand{\A}{\mathbb{A}}
\newcommand{\PP}{\mathbb{P}}
\newcommand{\E}{\mathbb{E}}
\newcommand{\eps}{\varepsilon}
\newcommand{\dd}{\,\mathrm{d}}
\begin{document}
\section*{UnsolvedMath ID: 6800003}
\textbf{Source problem code:} catalogue entry
\subsection{Definitions and mathematical statement}
Let $M$ be a connected closed flat Riemannian $n$-manifold. Its universal cover is Euclidean space and $M\cong\R^n/\Gamma$, where $\Gamma$ is a torsion-free discrete cocompact group of Euclidean isometries. Write each deck transformation as $x\mapsto A_\gamma x+b_\gamma$. The holonomy representation is the linear representation
\[
\rho:\Gamma\longrightarrow O(n),\qquad \rho(\gamma)=A_\gamma.
\]
Its image is finite, and reducibility means existence of a real linear subspace $V$ with
\[
0<\dim V<n,\qquad A_\gamma V=V\quad\text{for every }\gamma\in\Gamma.
\]
This is reducibility of the full holonomy representation, not just triviality of restricted holonomy for a flat connection.

The source asks for an alternative geometric proof of this known reducibility statement. It does not ask to determine whether the theorem is true or to construct a new manifold. The phrase geometric proof has no formal criterion here; no particular allowed or forbidden algebraic tool is prescribed.

The dimension restriction is missing in the source. With the usual proper-nonzero-subspace definition, the one-dimensional trivial representation is irreducible, so the statement needs $n\geq2$ or a special convention for that case. The supplied summary also suggests a center/centralizer argument without hypotheses guaranteeing a nonzero common fixed subspace. Reducibility does not require an invariant vector, and no such proof is inferred from that sketch.
\subsection{Short English statement}
Give a more geometric proof of the established reducibility of full holonomy for closed flat manifolds in the intended dimensions. Clarify the dimension-one exception and the desired meaning of geometric proof.
\subsection{Sources}
\begin{itemize}
\item Local UnsolvedMath registry, permanent ID 6800003, exact source code catalogue entry, ``Reducibility of the holonomy of flat manifolds''; proposed by Renato Bettiol, City University of New York (Lehman College), 2018. Registry status open, classification OPEN-TRIAGE; the supplied summary identifies the theorem itself as established and the requested proof style as the open task.
\item Morgan and Pansu, \emph{A List of Open Problems in Differential Geometry} (2018), question environment 3, \url{https://www.imo.universite-paris-saclay.fr/~pansu/problems_MTDG.tex}. Supplied source-TeX extraction accessed July 29, 2026.
\item The supplied summary points to Auslander--Kuranishi and Charlap, \emph{Bieberbach Groups and Flat Manifolds}, for Bieberbach-group treatments, without precise theorem, edition, or page citations. Its center/centralizer sketch does not provide a complete argument or its necessary hypotheses.
\item Provenance flags current status and redistribution terms NEEDS\_REVIEW; the source TeX is publicly available on a coauthor's publication page.
\end{itemize}
\end{document}
